% ===================================================================
%  ਸਾਰਣੀ ਨੰ. 8 — ਵਾਢੀ। — ਸ਼ਿਕਾਰ, ਅੰਗੂਰਾਂ ਦੀ ਵਾਢੀ, ਮੱਛੀ ਫੜਨਾ।
%  Traduction 2026 d'apres texto/io, controlee sur texto/fr, texto/en
%  et texto/hi. Consigne : voir texto/pa/10-sarni-01.tex.
%
%  L'appel de note est dans le TITRE de la scene — « ਵਾਢੀ (*) » — et
%  non dans un alinea : c'est ainsi que l'ido le compose.
% ===================================================================

%%K t08-noto-1 noto
\VUnotes{143.2mm}{%
\textit{(*) ਕਿਉਂਕਿ ਇਹ ਸ਼ਬਦ ਸਾਫ਼ ਨਹੀਂ: ਸਣ ਦੀ ਵਾਢੀ? ਅੰਗੂਰਾਂ ਦੀ? ਸੇਬਾਂ ਦੀ? ਸ਼ਾਇਦ
ਚੰਗਾ ਹੋਵੇ ਕਿ «} \VUgras{mesono} \textit{» ਕੰਮ ਵਿਚ ਲਿਆਂਦਾ ਜਾਵੇ, ਜਿਸ ਦੀ ਸਲਾਹ}
Mondo \textit{XIV, ਸਫ਼ਾ 242 ਵਿਚ ਦਿੱਤੀ ਗਈ ਹੈ। ਪਰ ਹੁਣ ਤਕ ਉਸ ਨਵੀਂ ਧਾਤੂ ਨੂੰ ਅਕਾਦਮੀ
ਨੇ ਅਪਣਾਇਆ ਨਹੀਂ ਤੇ ਉਹ ਇਦੋ ਲਹਿਰ ਦੇ ਆਮ ਨੇਮਾਂ ਮੁਤਾਬਕ ਬਹਿਸ ਦੀ ਉਡੀਕ ਵਿਚ ਹੈ।}%
}

%%K t08-apar-1 apar
\VUsaut{5.0mm}
\VUtitre{15.0pt}{60}{ਸਾਰਣੀ ਨੰ. 8}
\VUsaut{3.0mm}
\VUcentre{13.2pt}{90}{\textit{ਪਹਿਲਾ ਦ੍ਰਿਸ਼।}}
\VUsaut{3.0mm}
\VUpk{10.2pt}{40}{ਵਾਢੀ (*)।}
\VUsaut{4.0mm}
\VUpk{10.2pt}{40}{ਪਿੰਡ ਦਾ ਰੂਪ।}
\VUsaut{3.0mm}

%%K t08-c1-01-1 p
1. --- \VUgras{ਗਰਮੀ} ਨਾਲ ਪਿੰਡ ਨੇ ਫਿਰ ਇਕ ਵਾਰ ਆਪਣਾ ਰੂਪ ਬਦਲ ਲਿਆ ਹੈ। ਸਿਆੜ ਨੂੰ
ਸੌਂਪਿਆ ਬੀ ਫੁੱਟ ਪਿਆ; ਜ਼ਮੀਨ ਵਿਚੋਂ ਇਕ ਨਿੱਕਾ ਹਰਾ ਬੂਟਾ ਨਿਕਲਿਆ ਤੇ ਉਸ ਨੇ ਆਪਣਾ ਸਿੱਟਾ
ਦਿੱਤਾ। ਦਾਣਾ ਬਣਿਆ, ਫਿਰ ਪੱਕਿਆ, ਤੇ ਹੁਣ ਵਾਢੀ ਤੋਂ ਬਿਨਾਂ ਕੁਝ ਬਾਕੀ ਨਹੀਂ।

%%K t08-c1-02-1 p
2. --- \VUgras{ਜੰਗਲੀ ਫੁੱਲ} ਖਿੜ ਆਏ ਹਨ। \VUgras{ਨੀਲ-ਫੁੱਲ} \textsuperscript{(1)},
\VUgras{ਪੋਸਤ ਦੇ ਫੁੱਲ} \textsuperscript{(2)} ਤੇ \VUgras{ਡਿਜਿਟੈਲਿਸ}
\textsuperscript{(3)} ਕਣਕ ਦੇ ਖੇਤਾਂ ਦੇ ਇਕੋ ਜਿਹੇ ਰੰਗ ਉੱਤੇ ਆਪਣੇ ਤਾਜ਼ੇ \VUgras{ਫੁੱਲਾਂ}
ਦਾ ਖ਼ੁਸ਼ ਵਿਰੋਧ ਰੱਖਦੇ ਹਨ।

%%K t08-c1-03-1 p
3. --- ਫਲਾਂ ਵਾਲੇ ਰੁੱਖਾਂ ਨੇ ਪੱਤੇ ਕੱਢ ਲਏ ਹਨ; ਫਲ ਪੱਕ ਗਏ ਹਨ। ਨਿੱਕਾ
\VUgras{ਚੈਰੀ ਦਾ ਰੁੱਖ} \textsuperscript{(4)} ਸੋਹਣੀਆਂ \VUgras{ਚੈਰੀਆਂ}
\textsuperscript{(5)} ਨਾਲ ਲੱਦਿਆ ਹੈ, ਜਿਨ੍ਹਾਂ ਨੂੰ ਬੱਚੇ ਤੋੜਦੇ ਤੇ ਚਾਅ ਨਾਲ ਖਾਂਦੇ ਹਨ।

%%K t08-c1-04-1 p
4. --- ਇੱਜੜ ਹੁਣ ਸਾਰਾ ਦਿਨ ਬਾਹਰ ਲੰਘਾ ਸਕਦਾ ਹੈ।

%%K t08-c1-04-2 p
\VUgras{ਲੇਲੇ} \textsuperscript{(6)} ਵੱਡੇ ਹੋ ਕੇ ਸੋਹਣੀਆਂ \VUgras{ਭੇਡਾਂ}
\textsuperscript{(7)} ਤੇ ਸੋਹਣੇ \VUgras{ਛੱਤਰੇ} \textsuperscript{(8)} ਬਣ ਗਏ ਹਨ, ਸੰਘਣੀ
ਉੱਨ ਵਾਲੇ। ਉਹ \VUgras{ਚਰਾਗਾਹ} \textsuperscript{(9)} ਦਾ \VUgras{ਘਾਹ} ਸ਼ਾਂਤੀ ਨਾਲ ਚਰਦੇ
ਹਨ, ਜਿਸ ਉੱਤੇ \VUgras{ਗੁਲਬਹਾਰ} ਖਿੱਲਰੇ ਹਨ।

%%K t08-c1-04-3 p
ਛੇਤੀ ਹੀ ਇਨ੍ਹਾਂ ਡੰਗਰਾਂ ਦੀ \VUgras{ਉੱਨ} ਲਾਹੀ ਜਾਵੇਗੀ, ਜਿਸ ਤੋਂ ਸਾਡੇ ਕੱਪੜਿਆਂ ਦਾ ਕੱਪੜਾ
ਬਣਦਾ ਹੈ।

%%K t08-tit-2 sub
\VUsaut{5.0mm}
\VUpk{10.2pt}{40}{ਪਾਲੀ।}
\VUsaut{3.4mm}

%%K t08-c1-05-1 p
5. --- \VUgras{ਪਾਲੀ} \textsuperscript{(10)} ਉਨ੍ਹਾਂ ਵੱਲ ਵੱਧ ਧਿਆਨ ਦਿੰਦਾ ਨਹੀਂ ਲਗਦਾ।
ਉਸ ਦੀ ਇੱਕੋ ਫ਼ਿਕਰ ਹੈ ਦਿਸਹੱਦੇ ਵੱਲ ਜਾਂਦਾ ਤੂਫ਼ਾਨ, ਤੇ \VUgras{ਆਜੜੀ}
\textsuperscript{(13)}, ਜੋ ਆਪਣੀ \VUgras{ਸੁਰਾਹੀ} \textsuperscript{(14)} ਬੁੱਲ੍ਹਾਂ ਤਕ
ਚੁੱਕੀ ਹੋਇਆ ਹੈ, ਉਸ ਤੋਂ ਵੀ ਵੱਧ ਬੇਫ਼ਿਕਰ ਲਗਦਾ ਹੈ।

%%K t08-c1-06-1 p
6. --- ਗਰਮੀ ਸਾਹ ਘੁੱਟਣ ਵਾਲੀ ਹੈ; ਕਿਉਂਕਿ \VUgras{ਕੁੱਤਾ} \textsuperscript{(15)} ਵੀ
ਠੰਢਾ ਹੋਣ ਆਉਂਦਾ ਹੈ; ਉਹ \VUgras{ਨਾਲੇ} \textsuperscript{(16)} ਦਾ ਤਾਜ਼ਾ ਪਾਣੀ ਚੱਟਦਾ ਹੈ
ਤੇ ਇਕ ਵਿਚਾਰੇ ਨਿੱਕੇ \VUgras{ਡੱਡੂ} \textsuperscript{(17)} ਨੂੰ ਚੌਂਕਾ ਦਿੰਦਾ ਹੈ ਜੋ ਕੰਢੇ ਦੇ
ਘਾਹ ਵਿਚ ਦੁਬਕਿਆ ਹੋਇਆ ਸੀ। ਉਹ ਉਸ ਸਾਫ਼ ਪਾਣੀ ਵਿਚ ਟੱਪ ਪੈਂਦਾ ਹੈ ਜਿੱਥੇ \VUgras{ਕੰਵਲ}
\textsuperscript{(18)} ਖਿੜੇ ਹਨ।

%%K t08-c1-07-1 p
7. --- ਇਕ \VUgras{ਮਮੋਲਾ} \textsuperscript{(19)} ਉਸ ਨਿੱਕੇ \VUgras{ਚਸ਼ਮੇ}
\textsuperscript{(20)} ਕੋਲ ਨ੍ਹਾ ਰਿਹਾ ਹੈ ਜੋ ਦੋ ਚਟਾਨਾਂ ਵਿਚੋਂ ਫੁੱਟਦਾ ਹੈ ਤੇ
\VUgras{ਕੰਡਿਆਂ ਵਾਲੀਆਂ ਝਾੜੀਆਂ} \textsuperscript{(21)} ਤੇ \VUgras{ਬੇਰੀਆਂ ਦੀਆਂ
ਝਾੜੀਆਂ} \textsuperscript{(22)} ਦੇ ਹੇਠਾਂ ਲੁਕ ਜਾਂਦਾ ਹੈ।

%%K t08-tit-3 sub
\VUsaut{5.0mm}
\VUpk{10.2pt}{40}{ਵਾਢੀ ਦਾ ਕੰਮ।}
\VUsaut{3.4mm}

%%K t08-c1-08-1 p
8. --- ਪਿੰਡ ਦੇ ਬੱਚਿਆਂ ਲਈ ਕਣਕ ਦੀ ਵਾਢੀ ਵੱਡੇ ਪਰਚਾਵੇ ਦਾ ਵੇਲਾ ਹੈ। ਉਹ ਖ਼ੁਸ਼
\VUgras{ਕਲਾਬਾਜ਼ੀਆਂ} \textsuperscript{(24)} ਖਾਂਦੇ ਹਨ \VUgras{ਪਰਾਲੀ ਦੇ ਢੇਰਾਂ}
\textsuperscript{(23)} ਉੱਤੇ, \VUgras{ਵੱਢਣ ਵਾਲਿਆਂ} \textsuperscript{(26)} ਦੀ ਮਦਦ
ਕਰਦੇ ਹਨ, ਤੇ ਜਦੋਂ ਉਹ \VUgras{ਕਣਕ ਦਾ ਖੇਤ} \textsuperscript{(27)} ਵੱਢਦੇ ਹਨ ਆਪਣੀਆਂ
\VUgras{ਦਾਤਰੀਆਂ} \textsuperscript{(25)} ਨਾਲ, ਜੋ \VUgras{ਸਾਣ} \textsuperscript{(30)}
ਉੱਤੇ ਖ਼ੂਬ ਤੇਜ਼ ਕੀਤੀਆਂ ਗਈਆਂ ਹਨ, ਤਦੋਂ ਬੱਚੇ ਕਣਕ ਇਕੱਠੀ ਕਰ ਕੇ \VUgras{ਪੂਲੇ}
\textsuperscript{(31)} ਜਾਂ \VUgras{ਨਿੱਕੀਆਂ ਗੱਠੜੀਆਂ} \textsuperscript{(32)} ਬੰਨ੍ਹਦੇ
ਹਨ।

%%K t08-c1-09-1 p
9. --- ਕਦੇ ਕਦੇ ਉਨ੍ਹਾਂ ਵਿਚੋਂ ਵੱਡਿਆਂ ਨੂੰ \VUgras{ਦਾਤ} \textsuperscript{(12)} ਨਾਲ ਉਹ
\VUgras{ਸੁਨਹਿਰੀ ਡੰਠਲ} \textsuperscript{(28)} ਵੱਢਣ ਦਿੱਤੇ ਜਾਂਦੇ ਹਨ ਜੋ ਦਾਣਿਆਂ ਨਾਲ ਭਰੇ
ਭਾਰੇ \VUgras{ਸਿੱਟੇ} \textsuperscript{(29)} ਚੁੱਕੀ ਹੋਏ ਹਨ; ਤੇ ਨਿੱਕੇ,
\VUgras{ਡੰਗਰਾਂ} \textsuperscript{(33)} ਦੀ ਰਾਖੀ ਕਰਦੇ ਹੋਏ, ਜੋ \VUgras{ਮੈਦਾਨ}
\textsuperscript{(34)} ਵਿਚ ਵੱਡੇ \VUgras{ਲਿੰਡਨ ਦੇ ਰੁੱਖ} \textsuperscript{(35)} ਦੀ ਛਾਂ
ਵਿਚ ਚਰਦੇ ਹਨ, ਆਪਣੀ \VUgras{ਜਾਲੀ} \textsuperscript{(36)} ਵਿਚ ਸੋਹਣੀਆਂ
\VUgras{ਤਿੱਤਲੀਆਂ} ਫੜਦੇ ਹਨ।

%%K t08-c1-09-2 p
ਕੁਝ ਤਾਂ \VUgras{ਗੱਡੀ} \textsuperscript{(37)} ਉੱਤੇ ਚੜ੍ਹ ਕੇ ਉਹ ਪੂਲੇ ਜਮਾਉਂਦੇ ਹਨ ਜੋ
ਉਨ੍ਹਾਂ ਨੂੰ \VUgras{ਪੰਜੇ} \textsuperscript{(38)} ਉੱਤੇ ਚੁੱਕ ਕੇ ਦਿੱਤੇ ਜਾਂਦੇ ਹਨ।

%%K t08-c1-10-1 p
10. --- ਸਾਰੇ \VUgras{ਮੈਦਾਨ} \textsuperscript{(39)} ਵਿਚ, ਜਿੱਥੋਂ \VUgras{ਚੰਡੋਲਾਂ}
\textsuperscript{(41)} ਉੱਡਦੀਆਂ ਹਨ, ਵੱਡੀ ਚਹਿਲ-ਪਹਿਲ ਹੈ। ਸਾਰੇ ਵਾਢੀ ਵਿਚ ਲੱਗੇ ਹਨ। ਪਰ
ਜਿੱਥੇ ਗ਼ਰੀਬ ਪੁਰਾਣੇ ਔਜ਼ਾਰ ਹੀ ਕੰਮ ਵਿਚ ਲਿਆਉਂਦੇ ਹਨ — \VUgras{ਦਾਤਰੀਆਂ, ਦਾਤ} ਤੇ
\VUgras{ਮੁੰਗਲੀ} \textsuperscript{(40)} —, ਉੱਥੇ ਧਨੀ ਜ਼ਿਮੀਂਦਾਰ ਆਪਣੀ ਕਣਕ ਜਾਂ ਆਪਣੀ
\VUgras{ਰਾਈ} \textsuperscript{(43)} \VUgras{ਵੱਢਣ ਦੀ ਮਸ਼ੀਨ} \textsuperscript{(42)} ਨਾਲ
ਵਢਵਾਉਂਦਾ ਹੈ। ਫਿਰ, ਪੂਲਿਆਂ ਨੂੰ ਕੋਠੇ ਤਕ ਢੋਏ ਬਿਨਾਂ, ਉੱਥੇ ਹੀ ਵੱਡੇ \VUgras{ਢੇਰ}
\textsuperscript{(44)} ਲਾ ਦਿੱਤੇ ਜਾਂਦੇ ਹਨ।

%%K t08-c1-11-1 p
11. --- ਫਿਰ ਇਕ \VUgras{ਗਾਹੁਣ ਦੀ ਮਸ਼ੀਨ} \textsuperscript{(47)} ਲਿਆਂਦੀ ਜਾਂਦੀ ਹੈ,
ਜਿਸ ਨੂੰ ਇਕ \VUgras{ਚਲਦਾ ਇੰਜਣ} \textsuperscript{(45)} ਇਕ \VUgras{ਪੱਟੇ}
\textsuperscript{(46)} ਦੇ ਜ਼ਰੀਏ ਚਲਾਉਂਦਾ ਹੈ, ਤੇ ਇਸ ਤਰ੍ਹਾਂ ਦਾਣਾ \VUgras{ਪਰਾਲੀ} ਤੋਂ
ਵੱਖ ਹੋ ਜਾਂਦਾ ਹੈ, ਉਸੇ ਖੇਤ ਵਿਚ ਜਿੱਥੇ ਉਹ ਬੀਜਿਆ ਗਿਆ ਸੀ।

%%K t08-tit-4 sub
\VUsaut{5.0mm}
\VUpk{10.2pt}{40}{ਤੂਫ਼ਾਨ।}
\VUsaut{3.4mm}

%%K t08-c1-12-1 p
12. --- ਵਾਢੀ ਦੇ ਵੇਲੇ ਅਕਸਰ ਤਿੱਖੇ \VUgras{ਤੂਫ਼ਾਨ} \textsuperscript{(53)} ਟੁੱਟ ਪੈਂਦੇ
ਹਨ। ਪਹਿਲਾਂ ਦੂਰੋਂ \VUgras{ਬੱਦਲਾਂ ਦੀ ਗਰਜ} ਸੁਣਾਈ ਦਿੰਦੀ ਹੈ; ਇਕ ਤੇਜ਼
\VUgras{ਲਹਿੰਦੀ ਹਵਾ} \textsuperscript{(54)} ਉੱਠਦੀ ਹੈ ਤੇ \VUgras{ਜੰਗਲ}
\textsuperscript{(49)} ਦੇ ਸਭ ਤੋਂ ਵੱਡੇ ਰੁੱਖਾਂ ਨੂੰ ਕਰਾਹੁੰਦੇ ਹੋਏ ਝੁਕਾ ਦਿੰਦੀ ਹੈ। ਕਾਲੇ
\VUgras{ਬੱਦਲ} \textsuperscript{(56)} ਅਸਮਾਨ ਉੱਤੇ ਚੜ੍ਹ ਆਉਂਦੇ ਹਨ; ਉਨ੍ਹਾਂ ਨੂੰ
\VUgras{ਬਿਜਲੀ} \textsuperscript{(55)} ਦੀਆਂ ਚੁੰਧਿਆ ਦੇਣ ਵਾਲੀਆਂ ਵਿੰਗੀਆਂ ਲਕੀਰਾਂ ਚੀਰਦੀਆਂ
ਹਨ। \VUgras{ਮੀਂਹ} ਤੇ ਕਦੇ \VUgras{ਗੜੇ} ਪੈਣ ਲੱਗ ਪੈਂਦੇ ਹਨ, ਤੇ ਵੱਢਣ ਨੂੰ ਤਿਆਰ ਫ਼ਸਲ
ਨੂੰ ਵਿਛਾ ਦਿੰਦੇ ਹਨ।

%%K t08-c1-13-1 p
13. --- ਅਕਸਰ ਬਿਜਲੀ ਕਿਸੇ ਇਮਾਰਤ ਨੂੰ ਅੱਗ ਲਾ ਦਿੰਦੀ ਹੈ। ਪਿਛਲੇ ਵਰ੍ਹੇ, ਮਿਸਾਲ ਲਈ,
\VUgras{ਪੌਣ-ਚੱਕੀ} \textsuperscript{(50)} ਦੀ ਛੱਤ ਸੜ ਕੇ ਸੁਆਹ ਹੋ ਗਈ ਤੇ ਉਸ ਦੇ ਦੋ
\VUgras{ਖੰਭ} \textsuperscript{(51)} ਚੂਰ-ਚੂਰ ਹੋ ਗਏ।

%%K t08-c1-14-1 p
14. --- ਕੁਝ ਘੰਟਿਆਂ ਪਿੱਛੋਂ ਸ਼ਾਂਤੀ ਮੁੜ ਆਉਂਦੀ ਹੈ। ਦਿਸਹੱਦੇ ਉੱਤੇ ਇਕ ਚਮਕਦੀ
\VUgras{ਪੀਂਘ} \textsuperscript{(52)} ਵਿਖਾਈ ਦਿੰਦੀ ਹੈ, ਤੇ ਕੁਦਰਤ ਉਹ ਜੀਵਨ ਮੁੜ ਚੁੱਕ
ਲੈਂਦੀ ਹੈ ਜੋ ਥੋੜੀ ਦੇਰ ਲਈ ਰੁਕ ਗਿਆ ਸੀ। ਜੋ ਪਿੰਡ ਵਾਲੇ \VUgras{ਝੁੱਗੀਆਂ}
\textsuperscript{(48)} ਵਿਚ ਆ ਲੁਕੇ ਸਨ, ਉਹ ਕੰਮ ਉੱਤੇ ਮੁੜਦੇ ਹਨ ਤੇ ਹੁਣੇ ਹੁਣੇ ਹੋਏ ਨੁਕਸਾਨ
ਦੀ ਮੁਰੰਮਤ ਵਿਚ ਹਿੰਮਤ ਨਾਲ ਜੁਟ ਜਾਂਦੇ ਹਨ।

%%K t08-tit-5 sub
\VUsaut{6.0mm}
\VUcentre{13.2pt}{90}{\textit{ਦੂਜਾ ਦ੍ਰਿਸ਼।}}
\VUsaut{3.6mm}

%%K t08-tit-6 sub
\VUpk{10.2pt}{40}{ਸ਼ਿਕਾਰ। --- ਅੰਗੂਰਾਂ ਦੀ ਵਾਢੀ।}
\VUsaut{1.4mm}
\VUpk{10.2pt}{40}{ਮੱਛੀ ਫੜਨਾ।}
\VUsaut{4.0mm}

%%K t08-c2-01-1 p
1. --- \VUgras{ਅਗਸਤ} ਦੇ ਅਖ਼ੀਰ ਜਾਂ \VUgras{ਸਤੰਬਰ} ਦੇ ਵਿਚਕਾਰ, ਇਲਾਕੇ ਮੁਤਾਬਕ, ਸ਼ਿਕਾਰ
ਦਾ ਮੌਸਮ ਖੁੱਲ੍ਹਦਾ ਹੈ। ਸੂਰਜ ਦੀਆਂ ਪਹਿਲੀਆਂ ਕਿਰਨਾਂ \VUgras{ਕਿਰਲੀਆਂ}
\textsuperscript{(2)} ਨੂੰ ਉਨ੍ਹਾਂ ਦੀਆਂ ਖੁੱਡਾਂ ਤੋਂ ਬਾਹਰ ਲਿਆਉਂਦੀਆਂ ਵੀ ਨਹੀਂ ਕਿ
\VUgras{ਸ਼ਿਕਾਰੀ} \textsuperscript{(5)} ਆਪਣੇ \VUgras{ਕੁੱਤਿਆਂ} \textsuperscript{(4)}
ਨਾਲ ਨਿਕਲ ਪੈਂਦਾ ਹੈ।

%%K t08-c2-02-1 p
2. --- ਉਸ ਨੇ ਮੋਟੇ ਕੱਪੜੇ ਦਾ ਆਪਣਾ ਤਕੜਾ \VUgras{ਸ਼ਿਕਾਰੀ ਸੂਟ} \textsuperscript{(9)}
ਪਾਇਆ ਹੈ, ਤੇ ਚਮੜੇ ਦੇ \VUgras{ਪੱਟੇ}, ਜੋ ਉਸ ਨੂੰ ਉਸ ਗਿੱਲੇ ਘਾਹ ਵਿਚ ਤੁਰਨ ਦੇਣਗੇ ਜਿੱਥੇ
\VUgras{ਡੱਡੂ} \textsuperscript{(1)} ਲੁਕੇ ਰਹਿੰਦੇ ਹਨ; ਉਸ ਨੇ ਆਪਣੀ \VUgras{ਬੰਦੂਕ}
\textsuperscript{(6)} ਤੇ ਆਪਣਾ \VUgras{ਸ਼ਿਕਾਰ ਦਾ ਥੈਲਾ} \textsuperscript{(10)} ਚੁੱਕਿਆ
ਹੈ, ਤੇ ਤੁਰ ਪਿਆ ਹੈ।

%%K t08-c2-03-1 p
3. --- ਸ਼ਿਕਾਰ ਦੀ ਉਮੰਗ ਉਸ ਨੂੰ ਇਕ ਅੰਗੂਰਾਂ ਦੇ ਬਾਗ਼ ਵਿਚਕਾਰ ਲੈ ਆਉਂਦੀ ਹੈ ਜਿੱਥੇ ਅੰਗੂਰਾਂ
ਦੀ ਵਾਢੀ ਪੂਰੇ ਜ਼ੋਰ ਉੱਤੇ ਹੈ। ਅੰਗੂਰ ਉਗਾਉਣ ਵਾਲੇ ਦੇ ਬੱਚੇ, ਜੋ ਆਮ ਤੌਰ ਉੱਤੇ ਸੜਕ ਦੇ ਕੰਢੇ
ਬੱਕਰੀਆਂ ਚਾਰਦੇ ਹਨ, ਅੱਜ ਉਨ੍ਹਾਂ ਨੂੰ ਜਿੱਥੇ ਜੀਅ ਕਰੇ ਚਰਨ ਦੇ ਰਹੇ ਹਨ, ਤੇ ਇਕ ਬੁੱਢੇ
\VUgras{ਬੱਕਰੇ} \textsuperscript{(3)} ਨੂੰ ਕਿੱਲੇ ਨਾਲ ਬੰਨ੍ਹ ਕੇ — ਜੋ ਆਪਣੀ ਆਜ਼ਾਦੀ ਦਾ
ਫ਼ਾਇਦਾ ਚੁੱਕ ਕੇ ਜ਼ਰੂਰ ਨੱਸ ਜਾਂਦਾ — ਉਹ ਵੀ ਆਪਣੇ ਹਿੱਸੇ ਦਾ ਮਜ਼ਾ ਲੈ ਰਹੇ ਹਨ। ਇਕ
\VUgras{ਅੰਗੂਰਾਂ ਦੀ ਟੋਕਰੀ} \textsuperscript{(12)} ਤੋਂ ਇਕ ਗੁੱਛਾ ਲੈਂਦਾ ਹੈ ਤੇ ਉਸ ਦਾ
\VUgras{ਰਸ} \textsuperscript{(14)} ਇਕ \VUgras{ਪਿਆਲੇ} \textsuperscript{(13)} ਵਿਚ ਨਿਚੋੜ
ਕੇ ਮਸਤ (ਬਿਨਾਂ ਖ਼ਮੀਰ ਦੀ ਸ਼ਰਾਬ) ਬਣਾਉਣ ਵਿਚ ਦਿਲ ਪਰਚਾਉਂਦਾ ਹੈ। ਦੂਜਾ ਆਪਣੇ ਸਿਰ ਉੱਤੇ
\VUgras{ਪੱਤਿਆਂ ਦਾ ਤਾਜ} \textsuperscript{(15)} ਰੱਖ ਰਿਹਾ ਹੈ ਜੋ ਉਸ ਨੇ ਹੁਣੇ ਗੁੰਦਿਆ ਹੈ।
ਉਨ੍ਹਾਂ ਕੋਲ ਢੀਠ \VUgras{ਚਿੜੀਆਂ} \textsuperscript{(16)} ਅੰਗੂਰ ਚੋਰੀ ਕਰਨ ਕੋਹਲੂ ਤਕ ਚਲੀਆਂ
ਆਉਂਦੀਆਂ ਹਨ।

%%K t08-c2-04-1 p
4. --- ਜਦੋਂ ਤਕ ਬੱਚੇ ਦਿਲ ਪਰਚਾਉਂਦੇ ਹਨ, ਬੰਦੇ ਕੰਮ ਕਰਦੇ ਹਨ।
\VUgras{ਅੰਗੂਰ ਤੋੜਨ ਵਾਲੇ} \textsuperscript{(18)} ਅੰਗੂਰ \VUgras{ਪਿੱਠ ਦੀਆਂ ਟੋਕਰੀਆਂ}
\textsuperscript{(17)} ਵਿਚ ਤਹਿਖ਼ਾਨੇ ਤਕ ਲੈ ਜਾਂਦੇ ਹਨ; ਇਕ \VUgras{ਪੀਪਾ-ਸਾਜ਼}
\textsuperscript{(19)} \VUgras{ਪੀਪਿਆਂ} \textsuperscript{(20)} ਦੀ ਮੁਰੰਮਤ ਕਰਦਾ ਹੈ,
ਵੇਖਦਾ ਹੈ ਕਿ \VUgras{ਫੱਟੇ} \textsuperscript{(22)} ਤੇ \VUgras{ਢੱਕਣ}
\textsuperscript{(21)} ਠੀਕ ਹਨ, ਤੇ ਟੁੱਟੇ \VUgras{ਛੱਲੇ} \textsuperscript{(23)} ਬਦਲਦਾ
ਹੈ। ਉਹ ਵੱਡੇ \VUgras{ਕੁੰਡ} \textsuperscript{(25)} ਵੀ ਉਸੇ ਨੇ ਬਣਾਏ ਸਨ, ਜਿਨ੍ਹਾਂ ਵਿਚ
ਖ਼ਮੀਰ ਉਠਾਉਣ ਨੂੰ \VUgras{ਅੰਗੂਰ} \textsuperscript{(26)} ਪਾਏ ਜਾਂਦੇ ਹਨ।

%%K t08-c2-05-1 p
5. --- ਖ਼ਮੀਰ ਉੱਠ ਚੁੱਕਣ ਉੱਤੇ ਸ਼ਰਾਬ ਕੱਢ ਲਈ ਜਾਂਦੀ ਹੈ ਤੇ ਬਚਿਆ ਹੋਇਆ \VUgras{ਕੋਹਲੂ}
\textsuperscript{(28)} ਉੱਤੇ ਰੱਖਿਆ ਜਾਂਦਾ ਹੈ; ਵੱਡੇ \VUgras{ਪੇਚ} \textsuperscript{(29)}
ਨੂੰ ਹੌਲੀ ਹੌਲੀ ਕੱਸ ਕੇ ਰਸ ਨਿਚੋੜਿਆ ਜਾਂਦਾ ਹੈ, ਜੋ ਇਕ \VUgras{ਟੱਬ}
\textsuperscript{(32)} ਵਿਚ ਵਗਦਾ ਹੈ ਤੇ ਹੋਰ ਵੀ \VUgras{ਸ਼ਰਾਬ} \textsuperscript{(31)}
ਦਿੰਦਾ ਹੈ।

%%K t08-tit-8 sub
\VUsaut{6.0mm}
\VUpk{10.2pt}{40}{ਬੀਅਰ ਤੇ ਸੇਬ ਦੀ ਸ਼ਰਾਬ।}
\VUsaut{3.4mm}

%%K t08-c2-06-1 p
6. --- \VUgras{ਤਹਿਖ਼ਾਨੇ} \textsuperscript{(27)} ਦੀ ਕੰਧ ਉੱਤੇ \VUgras{ਹੌਪਸ}
\textsuperscript{(30)} ਦੀ ਇਕ ਟਾਹਣੀ ਚੜ੍ਹਦੀ ਹੈ। ਇੱਥੇ ਉਹ ਸਿਰਫ਼ ਸਜਾਵਟੀ ਬੂਟਾ ਹੈ, ਪਰ ਕੁਝ
ਦੇਸਾਂ ਵਿਚ, ਜਿੱਥੇ ਵੇਲ ਬਹੁਤ ਘੱਟ ਹੁੰਦੀ ਹੈ, \VUgras{ਬੀਅਰ} ਬਣਾਉਣ ਲਈ ਹੌਪਸ ਉਗਾਏ ਜਾਂਦੇ
ਹਨ।

%%K t08-c2-07-1 p
7. --- ਨਿੱਕਾ \VUgras{ਸੇਬ ਦਾ ਰੁੱਖ} \textsuperscript{(33)} ਸੇਬਾਂ ਨਾਲ ਲੱਦਿਆ ਹੈ,
ਜਿਨ੍ਹਾਂ ਤੋਂ ਪੱਕਣ ਉੱਤੇ ਵਧੀਆ \VUgras{ਸੇਬ ਦੀ ਸ਼ਰਾਬ} ਬਣੇਗੀ। ਆਪਣੇ \VUgras{ਪਿੰਜਰੇ}
\textsuperscript{(34)} ਵਿਚ \VUgras{ਕਨੇਰੀ} \textsuperscript{(35)} ਤੇ
\VUgras{ਗੋਲਡਫ਼ਿੰਚ} \textsuperscript{(36)} ਉਨ੍ਹਾਂ ਚਿੜੀਆਂ ਨੂੰ ਸੜ ਨਾਲ ਵੇਖਦੇ ਹਨ ਜੋ
ਆਜ਼ਾਦ ਟੱਪੂਸੀਆਂ ਮਾਰ ਰਹੀਆਂ ਹਨ।

%%K t08-c2-08-1 p
8. --- ਜਿੱਥੋਂ ਤਕ ਨਜ਼ਰ ਜਾਂਦੀ ਹੈ, ਟਿੱਬੀਆਂ ਦੀ ਢਲਾਣ ਉੱਤੇ, \VUgras{ਵੇਲ ਦੇ ਕਿੱਲਿਆਂ}
\textsuperscript{(40)} ਦੀਆਂ ਕਤਾਰਾਂ ਖੜੀਆਂ ਹਨ, ਜੋ ਸ਼ਾਨਦਾਰ \VUgras{ਵੇਲ ਦੇ ਮੁੱਢਾਂ}
\textsuperscript{(39)} ਨੂੰ ਸੰਭਾਲੀ ਬੈਠੀਆਂ ਹਨ, \VUgras{ਗੁੱਛਿਆਂ} ਨਾਲ ਲੱਦੇ ਹੋਏ।

%%K t08-c2-08-2 p
ਸਾਰਾ ਵਰ੍ਹਾ \VUgras{ਅੰਗੂਰ ਉਗਾਉਣ ਵਾਲੇ} \textsuperscript{(42)} ਨੇ ਆਪਣੇ
\VUgras{ਅੰਗੂਰਾਂ ਦੇ ਬਾਗ਼} \textsuperscript{(38)} ਦੀ ਸੰਭਾਲ ਵਿਚ ਮਿਹਨਤ ਕੀਤੀ ਹੈ, ਤੇ ਹੁਣ
ਉਸ ਨੂੰ ਚੰਗੀ ਫ਼ਸਲ ਤੋਂ ਉਸ ਦਾ ਇਨਾਮ ਮਿਲਦਾ ਹੈ। \VUgras{ਤੋੜਨ ਵਾਲੀਆਂ}
\textsuperscript{(41)} ਉਹ ਕੁੰਡ ਭਰਦੀਆਂ ਹਨ ਜੋ \VUgras{ਖੱਚਰਾਂ} \textsuperscript{(43)}
ਨਾਲ ਖਿੱਚੀ ਗੱਡੀ ਉੱਤੇ ਰੱਖੇ ਹਨ, ਤੇ ਸਾਰੇ ਖ਼ੁਸ਼ ਹਨ।

%%K t08-tit-9 sub
\VUsaut{5.0mm}
\VUpk{10.2pt}{40}{ਮੱਛੀ ਫੜਨਾ।}
\VUsaut{3.4mm}

%%K t08-c2-09-1 p
9. --- ਇਹੋ ਗੱਲ \VUgras{ਬੰਸੀ ਵਾਲਿਆਂ} \textsuperscript{(44)} ਦੀ ਹੈ, ਜੋ ਪਹੁ ਫੁੱਟਦੇ ਹੀ
ਆਪਣੀਆਂ \VUgras{ਬੰਸੀਆਂ} \textsuperscript{(45)} ਨਾਲ ਪਾਣੀ ਦੇ ਕੰਢੇ ਆ ਜੰਮੇ ਹਨ।

%%K t08-c2-09-2 p
\VUgras{ਮੱਛੀਆਂ} \textsuperscript{(47)} ਕੁੰਡਾ ਫੜ ਲੈਂਦੀਆਂ ਹਨ ਜਿਉਂ ਹੀ ਉਹ ਆਪਣੀ
\VUgras{ਡੋਰ} \textsuperscript{(46)} ਪਾਣੀ ਵਿਚ ਪਾਉਂਦੇ ਹਨ, ਤੇ ਉਨ੍ਹਾਂ ਨੂੰ
\VUgras{ਚੋਗਾ} ਬਦਲਣ ਦਾ ਵੇਲਾ ਵੀ ਨਹੀਂ ਮਿਲਦਾ ਕਿ ਡੋਰ ਦੇ ਸਿਰੇ ਉੱਤੇ ਇਕ ਨਵਾਂ ਸ਼ਿਕਾਰ ਤੜਪ
ਰਿਹਾ ਹੁੰਦਾ ਹੈ।

%%K t08-c2-09-3 p
ਫਿਰ ਵੀ \VUgras{ਜਾਲ ਵਾਲਾ ਮਾਹੀਗੀਰ} \textsuperscript{(48)}, ਜੋ ਆਪਣੀ \VUgras{ਬੇੜੀ}
\textsuperscript{(49)} ਤੋਂ \VUgras{ਦਰਿਆ} \textsuperscript{(50)} ਦੇ ਵਿਚਕਾਰ ਜਾਲ ਸੁੱਟਦਾ
ਹੈ, ਕਿਤੇ ਵੱਧ ਫੜਦਾ ਹੈ।

%%K t08-c2-10-1 p
10. --- ਪੱਤਝੜ ਚੰਗੀ ਸੈਰ ਦਾ ਮੌਸਮ ਹੈ: ਪੈਦਲ, \VUgras{ਘੋੜੇ} \textsuperscript{(52)}
ਉੱਤੇ, \VUgras{ਸਾਈਕਲ} \textsuperscript{(53)} ਨਾਲ, \VUgras{ਮੋਟਰ} \textsuperscript{(54)}
ਨਾਲ। \VUgras{ਸੜਕਾਂ} \textsuperscript{(51)} ਸਾਫ਼ ਹਨ ਤੇ ਲੋਕ ਆਖ਼ਰੀ ਚੰਗੇ ਦਿਨਾਂ ਦਾ ਪੂਰਾ
ਫ਼ਾਇਦਾ ਚੁੱਕਦੇ ਹਨ, ਕਿਉਂਕਿ \VUgras{ਅਬਾਬੀਲਾਂ} \textsuperscript{(56)} ਪਹਿਲਾਂ ਹੀ ਛੱਤਾਂ
ਉੱਤੇ ਜੁੜਨ ਲੱਗ ਪਈਆਂ ਹਨ, ਨਿੱਘੇ ਦੇਸਾਂ ਨੂੰ ਉੱਡ ਜਾਣ ਲਈ। ਬੁੱਢੇ \VUgras{ਅਖ਼ਰੋਟ ਦੇ ਰੁੱਖ}
\textsuperscript{(57)} ਦੇ ਪੱਤੇ ਪੀਲੇ ਪੈ ਰਹੇ ਹਨ, \VUgras{ਅਖ਼ਰੋਟ} ਡਿੱਗ ਰਹੇ ਹਨ, ਤੇ ਇਕ
\VUgras{ਝੁੰਡ} \textsuperscript{(58)} ਵਿਚ ਖੜੇ ਉੱਚੇ \VUgras{ਰੁੱਖ}, ਜੋ
\VUgras{ਉੱਚੀ ਜ਼ਮੀਨ} \textsuperscript{(59)} ਉੱਤੇ ਹੈ, ਛੇਤੀ ਹੀ ਨੰਗੇ ਹੋ ਜਾਣਗੇ।

%%K t08-c2-11-1 p
11. --- \VUgras{ਸੂਰਜ} \textsuperscript{(60)} ਦੇਰ ਨਾਲ ਚੜ੍ਹਦਾ ਹੈ, ਦਿਨ ਨਿੱਕੇ ਹੁੰਦੇ
ਜਾਂਦੇ ਹਨ, \VUgras{ਸਵੇਰ} ਇਕ ਚੰਗੀ-ਖ਼ਾਸੀ ਠੰਢ ਮਹਿਸੂਸ ਹੋਣ ਲੱਗ ਪੈਂਦੀ ਹੈ, ਤੇ
\VUgras{ਚਾਹਾਂ} \textsuperscript{(61)} ਦੇ ਝੁੰਡ ਹੁਣੇ ਤੋਂ ਸਾਡਾ ਬੇਰਹਿਮ ਮੌਸਮ ਛੱਡ ਕੇ ਜਾ
ਰਹੇ ਹਨ।
\VUsaut{6.0mm}
\VUfilet{22mm}
